\section{Results}\label{sec:results}

Does training make the first proposed algorithm better on an unseen problem?
The available training trajectory in \cref{sec:method-status} records
validation behavior during learning. It does not answer the deployment
question: those validation outcomes inform admission and are not an
untouched test. This section specifies the comparisons still needed.
\TBD{Add independent evaluation evidence before making a performance claim.}
\TBD{Supply a source-bound benchmark before claiming deployment throughput.}

\paragraph{Earlier search diagnostics.}
The supplied method notes report 3,227 exact ties in 3,268 paired
policy-seeded versus generic final searches, plus a separate audit in which
98.84\% of reported winners were already in their starting seed populations.
The paired searches shared seeds, repeated development contexts and had
unequal wall times; they were not independent first-proposal draws. The two
historical inventories differ.

For a separate earlier model evaluated on development contexts, the notes
report fresh-program pairwise ranking accuracy of 0.506 with an interval
of [0.441, 0.568], which includes chance. The reference checkpoint and later adapter snapshot in
\cref{sec:method} require their own ranking evidence.
These observations describe historical search and ranking behavior; the
matched first-proposal test below is still required.
\TBD{Pin the inventories, checkpoints and seeding, timing and ranking
protocols, separately from the final comparisons below.}

\subsection{Evaluation contract}\label{sub:results-protocol}

Freeze the writer, permitted task inputs, output support, decoding rule and
budgets before opening final outcomes. Pin dataset versions and content
hashes, row splits and family assignments. Separate unseen problem
families from training and development, and report which task types were
actually tested.
Report known training and pretraining exposure for the writer, encoders
and learned comparators, with unresolved provenance marked explicitly.
Bound unseen-family claims to the exclusions actually established.
Each \EZ{} submission follows \cref{alg:gameday}; failed generations and
executions remain in the denominator. Report paired results and uncertainty
at the family level, alongside encoding, generation, fitting and prediction
costs. Disclose any internal model or hyperparameter selection in the
emitted procedure, even when the writer submits only one outer candidate.
\TBD{Document the final problem inventory, dataset identities, row splits,
family assignments and exclusions, component provenance, evaluator,
checkpoint, seeds, task metrics and failure accounting.}

\subsection{The effect of training}\label{sub:results-learning}

Test the learning claim with a frozen reference writer as the control.
Predeclare that checkpoint and the evaluated training intervention before
opening final outcomes. Hold the encoder, program grammar and output
support, permitted task inputs, decoding rules, fitting protocol, referee
and resource budgets constant; vary the writer checkpoint. Changes to the
support, encoder or evaluator belong in separate conditions.

Pair the first draws under predeclared seeds and report all of them, with
no selection of the best seed or candidate. Report paired family-level
outcomes, failures and costs alongside source-bound evidence of optimizer
updates and parameter changes. \TBD{Supply the matched first-proposal
comparison and training-update evidence before claiming transfer.}

\subsection{One proposal against fixed methods}\label{sub:results-headtohead}

The proposed comparison asks whether one task-conditioned program competes
with methods that run unchanged across tasks. Frontier-designed algorithms
must be fixed before final evaluation; record the originating models,
prompts, dates and implementation checks. Pin the KumoRFM configuration
\citep{kumorfm2025}, disclose its permitted inputs and fitting procedure, and
compare under the same task metrics. \Cref{tab:headtohead} reserves the
comparison; no winner is assumed.

\begin{table}[t]
  \centering
  \caption{\textbf{Planned paired comparisons.} No outcomes are yet supplied.
  Report ties, failures and excluded cases alongside wins.}
  \label{tab:headtohead}
  \small
  \begin{tabularx}{\linewidth}{@{}>{\raggedright\arraybackslash}X>{\raggedright\arraybackslash}X@{}}
    \toprule
    Comparator & Outcome and uncertainty \\
    \midrule
    Claude-designed method & \TBD{paired evaluation} \\
    Codex-designed method & \TBD{paired evaluation} \\
    Gemini-designed method & \TBD{paired evaluation} \\
    KumoRFM & \TBD{paired evaluation} \\
    \bottomrule
  \end{tabularx}
\end{table}

\begin{figure}[t]
  \centering
  \fbox{\parbox{0.88\linewidth}{\TBD{Leaderboard: supply the evaluated methods,
  paired outcomes, rating protocol and uncertainty.}}}
  \caption{\textbf{Leaderboard pending.} Elo is a summary of a specified
  comparison population, not a substitute for paired task results.}
  \label{fig:leaderboard}
\end{figure}

\subsection{The cost of more attempts}\label{sub:results-search}

Compare the single proposal with agents allowed successive attempts and
development feedback. Keep final outcomes unavailable to those agents.
Count every candidate attempt, including incomplete or invalid outputs
and failed executions; a repair that emits a revised candidate is another
attempt. Charge all task-time computation. Plot
performance against both attempts and cost; if an agent does not reach a
reference performance within its budget, report that limit.
\TBD{Supply attempt curves, matched budgets, repeated seeds and uncertainty.}

\begin{figure}[t]
  \centering
  \fbox{\parbox{0.88\linewidth}{\TBD{Search curve: performance versus attempts
  and compute, including failed attempts and unreached targets.}}}
  \caption{\textbf{Deployment search comparison pending.} The single-candidate
  condition fixes a budget; its measured quality remains to be established.}
  \label{fig:searchone}
\end{figure}

\subsection{What the map contributes}\label{sub:results-ablation}

A representation study should test whether nearby problem encodings share
useful programs, with data disjoint from fitting the map. Ranking agreement
should also be compared with the program-only head, so a global preference
for common programs is distinguishable from reading the problem.
\TBD{Supply problem-similarity and program-only-control evaluations.}

A ranking head can agree with a referee without improving the writer.
In the checkpoint described in \cref{sub:method-map}, the map neither
conditions nor screens the writer's draw. Its ranking accuracy therefore
cannot establish a reduction in wasted proposals. A causal ablation needs
a declared route by which the map changes training, with the writer,
training budget and deployment protocol otherwise matched.
\TBD{Identify an implemented influence path or limit the study to ranking.}

\begin{table}[t]
  \centering
  \caption{\textbf{Map evidence still required.} Ranking quality and a change
  in first-proposal quality are separate outcomes.}
  \label{tab:ablation}
  \small
  \begin{tabularx}{\linewidth}{@{}>{\raggedright\arraybackslash}X>{\raggedright\arraybackslash}X@{}}
    \toprule
    Question & Required evidence \\
    \midrule
    Does the map rank candidates? & \TBD{held-out ranking study} \\
    Does it improve the writer? & \TBD{matched causal ablation} \\
    Does it reduce execution failures? & \TBD{failure-class comparison} \\
    \bottomrule
  \end{tabularx}
\end{table}

\subsection{Failure cases and a worked example}\label{sub:results-failures}

Describe unsuccessful proposals and coverage limits as well as wins.
Assign each attempt one terminal status: incomplete generation, parse or
type-checking refusal, execution
or resource failure, invalid output under declared checks, and valid ties,
losses and wins. Every attempted candidate contributes its incurred cost.
Unresolved outcomes retain an explicit unknown status; fix their treatment
and the validity rules before opening final outcomes.
\TBD{Supply attempt inventories and costs by outcome class, identifying
unresolved cases.}

Choose the worked example by a rule fixed before inspection, and show its emitted
program, fitting procedure and comparison outcome; if the selected category
is empty, report that fact. \TBD{Supply the example selection rule, actual
program and verified outcomes, with no replacement of failed attempts.}
These results will determine which interpretations in \cref{sec:meaning}
the evidence can support.
